\documentclass[12pt, a4paper]{article}

% --- Packages ---
\usepackage[margin=1in]{geometry}
\usepackage{amsmath, amssymb}
\usepackage{graphicx}
\usepackage{hyperref}
\usepackage{cite}
\usepackage{float}
\usepackage{listings}
\usepackage{xcolor}
\usepackage{booktabs}
\usepackage{caption}
\usepackage{setspace}
\usepackage{authblk}
\usepackage{abstract}

\hypersetup{
    colorlinks=true,
    linkcolor=blue,
    citecolor=blue,
    urlcolor=blue
}

\lstset{
    language=Python,
    basicstyle=\ttfamily\small,
    keywordstyle=\color{blue},
    commentstyle=\color{gray},
    stringstyle=\color{red!60!black},
    numbers=left,
    numberstyle=\tiny\color{gray},
    breaklines=true,
    frame=single,
    captionpos=b
}

\onehalfspacing

% --- Title ---
\title{\textbf{Self-Organized Criticality in the Forest Fire Model:\\
A Complexity Science Perspective}}

\author{Adhya Makkar}
\affil[1]{Department of Chemical Engineering, IIT Delhi\\
Entry Number: 2021CH70028}

\date{\today}

% ==============================================================
\begin{document}

\maketitle
\thispagestyle{empty}

% --- Abstract ---
\begin{abstract}
\noindent
This paper applies the framework of complexity science, specifically the theory of
Self-Organized Criticality (SOC), to a cellular-automaton forest fire model. The
system is characterized by three fundamental SOC properties: noise (stochastic tree
growth and random lightning strikes), avalanches (cascading fire-spread events), and
connectivity (the density of trees governing fire propagation). A Python simulation
on a $50 \times 50$ periodic grid is used to study the frequency distribution of
avalanche sizes. The resulting log-log plot reveals a power-law-like distribution,
consistent with scale-free dynamics and SOC behavior. The relationship between tree
density (connectivity) and avalanche magnitude is analyzed, and the system's
tendency to self-organize toward a critical state without external tuning is
discussed. The results confirm that the forest fire model is a canonical example of
SOC and offers broad applicability to real-world systems ranging from ecosystems to
financial markets.
\end{abstract}

\vspace{0.5em}
\noindent\textbf{Keywords:} Self-Organized Criticality, Forest Fire Model, Avalanche
Distribution, Cellular Automaton, Complexity Science, Power Law

\newpage
\tableofcontents
\newpage

% ==============================================================
\section{Introduction}
\label{sec:intro}

Complex systems are ubiquitous in nature—from earthquakes and neural networks to
stock markets and ecosystems. A striking feature of many such systems is their
tendency to evolve spontaneously to a \emph{critical state}, exhibiting scale-free
behavior without any fine-tuning of parameters. This phenomenon, termed
\textbf{Self-Organized Criticality} (SOC), was introduced by Bak, Tang, and
Wiesenfeld in their seminal 1987 paper through the sandpile model
\cite{bak1987self}.

SOC systems are characterized by three hallmark features: (i)~\textbf{noise}—small,
continuous perturbations driving the system; (ii)~\textbf{avalanches}—discrete,
catastrophic relaxation events whose sizes follow a power-law distribution; and
(iii)~\textbf{connectivity}—the local interaction structure that determines how
perturbations propagate.

The \textbf{forest fire model}, proposed by Drossel and Schwabl \cite{drossel1992self},
is one of the most intuitive and widely studied SOC models. Trees grow on a grid with
probability $p$ per time step; a lightning strike with probability $f \ll p$ can
ignite a tree, initiating a cascade of fire that spreads to all connected trees. The
ratio $p/f$ governs whether the system reaches a SOC state.

This paper presents a Python-based simulation of the forest fire model and analyses
its emergent SOC behavior, with a focus on the avalanche size distribution and its
relationship to tree density (connectivity). Section~\ref{sec:why_complex} justifies
the model as a complexity-science candidate. Section~\ref{sec:model} describes the
system. Section~\ref{sec:simulation} details the simulation. Section~\ref{sec:results}
presents and analyzes the results. Section~\ref{sec:soc} comments on the SOC state,
and Section~\ref{sec:conclusion} concludes.

% ==============================================================
\section{Why the Forest Fire Model is a Candidate for Complexity Science}
\label{sec:why_complex}

The forest fire model satisfies all the conditions that qualify a system for
complexity-science analysis:

\begin{enumerate}
    \item \textbf{Many interacting agents.} Each cell on the grid (tree, empty, or
    burning) interacts with its four nearest neighbors. Collective dynamics emerge
    from purely local rules with no global controller.

    \item \textbf{Emergent macroscopic behavior.} Individual cell transitions are
    simple, yet the system produces large-scale, correlated fire events whose
    statistics cannot be deduced from single-cell dynamics alone.

    \item \textbf{Scale-free avalanches.} Fire events span all length scales—from a
    single burning tree to system-spanning conflagrations—with no characteristic
    size. This is the hallmark of criticality.

    \item \textbf{Absence of a tuning parameter.} Unlike conventional phase
    transitions, the forest fire model organizes itself to the critical point for a
    wide range of $p$ and $f$ values, provided $f \ll p$.

    \item \textbf{Broad applicability.} The model is a metaphor for any
    contagion-relaxation system: disease spread, information cascades in social
    networks, and even supply-chain disruptions.
\end{enumerate}

These properties make the forest fire model an excellent pedagogical and research
vehicle for complexity science.

% ==============================================================
\section{Description of the System}
\label{sec:model}

\subsection{Model Overview}

The model is defined on a two-dimensional $L \times L$ lattice with periodic
(toroidal) boundary conditions. Each cell $i$ at position $(r,c)$ can be in one of
three discrete states:

\begin{equation}
    s_{r,c} \in \{\text{EMPTY}=0,\ \text{TREE}=1,\ \text{FIRE}=2\}
\end{equation}

The update rules, applied synchronously at each time step, are:

\begin{enumerate}
    \item \textbf{Fire spreading:} A TREE cell becomes FIRE if at least one of its
    four von Neumann neighbors is currently FIRE, or if a lightning strike occurs
    with probability $f$.
    \item \textbf{Tree growth:} An EMPTY cell becomes a TREE with probability $p$.
    \item \textbf{Fire extinguishing:} A FIRE cell becomes EMPTY in the next step
    (fire burns for exactly one time step).
\end{enumerate}

\subsection{Noise}

In SOC terminology, \emph{noise} refers to the slow, stochastic driving force that
pushes the system toward the critical state. In the forest fire model, noise is
embodied by:

\begin{itemize}
    \item \textbf{Tree growth} at rate $p = 0.05$ per cell per step—a slow,
    diffuse input of ``fuel'' into the system.
    \item \textbf{Lightning strikes} at rate $f = 0.001$—rare, random sparks that
    trigger relaxation events.
\end{itemize}

The separation of timescales ($f \ll p$) ensures that the forest can build up
significant connectivity between lightning events, enabling large avalanches.

\subsection{Avalanche}

An \emph{avalanche} in this system is a fire event—the cascade of burning cells
triggered by a single lightning strike propagating through connected trees. The
\emph{avalanche size} $S$ is defined as the number of cells that transition to the
FIRE state in a single time step. Avalanche sizes range from 1 (a single isolated
tree struck by lightning) to $O(L^2)$ (a system-spanning fire). The distribution
of avalanche sizes is the primary observable used to diagnose SOC.

\subsection{Connectivity}

\emph{Connectivity} refers to the spatial density and clustering of trees. A cell
is connected to its fire source if a path of adjacent TREE cells links it to a
burning cell. High tree density implies large connected clusters and thus larger
potential avalanches. Low density means fires are quickly contained by firebreaks
(EMPTY cells). The critical state occurs at an intermediate tree density at which
the forest is on the percolation threshold—fires can span the system, but do not
immediately consume every tree.

\subsection{Parameters}

\begin{table}[H]
\centering
\caption{Simulation parameters}
\begin{tabular}{lll}
\toprule
\textbf{Parameter} & \textbf{Symbol} & \textbf{Value} \\
\midrule
Grid size           & $L$     & 50 \\
Growth probability  & $p$     & 0.05 \\
Lightning probability & $f$   & 0.001 \\
Time steps          & $T$     & 500 \\
Boundary conditions & ---     & Periodic (toroidal) \\
\bottomrule
\end{tabular}
\label{tab:params}
\end{table}

% ==============================================================
\section{Computer Simulation}
\label{sec:simulation}

The simulation was implemented in Python 3 using \texttt{NumPy} for array operations
and \texttt{Matplotlib} for visualization. The complete source code is available in
the project GitHub repository (see Section~\ref{sec:code_link}).

\subsection{Algorithm}

At each of the $T = 500$ time steps, the algorithm:
\begin{enumerate}
    \item Creates a copy of the current grid.
    \item Iterates over all cells and applies the three update rules (fire spread,
    growth, extinguishing) using the state of the \emph{current} grid (synchronous
    update).
    \item Records the total number of cells transitioning to FIRE as the avalanche
    size for that step (if non-zero).
    \item Replaces the current grid with the updated copy.
\end{enumerate}

\subsection{Avalanche Size Distribution}

After the simulation, avalanche sizes are collected and binned. A log-log histogram
is plotted: $\log S$ on the $x$-axis versus $\log N(S)$ on the $y$-axis, where
$N(S)$ is the number of time steps with avalanche size $S$. A linear trend in this
plot is evidence of a power-law:

\begin{equation}
    N(S) \sim S^{-\tau}
\end{equation}

where $\tau$ is the scaling exponent characteristic of the SOC universality class.

\subsection{Source Code}
\label{sec:code_link}

The Python simulation code is reproduced below for completeness. The GitHub
repository is available at:
\url{https://github.com/Adhya2009/CLL798_Project_2021CH70028}



% ==============================================================
\section{Results and Analysis}
\label{sec:results}

\subsection{Avalanche Size Distribution}

The frequency distribution of avalanche sizes is shown in Figure~\ref{fig:loglog}.


The log-log plot reveals an approximately linear decrease in $N(S)$ with increasing
$S$, consistent with a power-law distribution $N(S) \sim S^{-\tau}$. This is the
defining signature of a system at a critical state. Small fires are frequent, while
large, system-spanning fires are rare—yet all sizes occur, with no preferred
scale.

\begin{figure}
    \centering
    \includegraphics[width=0.75\linewidth]{cll798 project simualtion image.png}
    \caption{Log-log plot of avalanche size frequency distribution $N(S)$ vs.\
    avalanche size}
    \label{fig:loglog}
\end{figure}


\subsection{Relationship Between Avalanche Size and Connectivity}

Tree density (connectivity) directly controls the maximum achievable avalanche size.
In the present simulation, the ratio $p/f = 50$ is large enough that the forest
rebuilds to a high density between fires. This elevated connectivity means fires
routinely encounter dense clusters and produce large avalanches, pushing the right
tail of the distribution to larger values.

Qualitatively, one can identify three regimes:
\begin{itemize}
    \item \textbf{Low connectivity ($\rho_{\text{trees}} \ll \rho_c$):} The forest
    is sparse; fires self-extinguish quickly. Avalanche distribution is
    approximately exponential—no power law.
    \item \textbf{Critical connectivity ($\rho_{\text{trees}} \approx \rho_c$):}
    The forest is at the percolation threshold. Avalanches span all scales; power-law
    distribution emerges.
    \item \textbf{High connectivity ($\rho_{\text{trees}} \gg \rho_c$):}
    System-spanning fires dominate; distribution is cut off at system size $L^2$.
\end{itemize}

The SOC mechanism ensures that the system \emph{self-regulates} to the critical
regime: fires reduce density when it is too high, while slow tree growth replenishes
it when it is too low.

% ==============================================================
\section{Self-Organized Critical State}
\label{sec:soc}

\subsection{Evidence for SOC}

The forest fire model exhibits the three canonical signatures of SOC:

\begin{enumerate}
    \item \textbf{Power-law avalanche distribution.} The log-log linearity of
    $N(S)$ versus $S$ confirms scale-free behavior, the most direct empirical
    evidence for criticality.

    \item \textbf{Self-organization without tuning.} Unlike a conventional
    second-order phase transition, criticality here is not achieved by tuning
    a single control parameter to a precise value. For any $f \ll p$, the system
    organizes itself to a statistically stationary critical state.

    \item \textbf{Long-range spatial correlations.} At the critical state, the
    size of connected tree clusters diverges (in the thermodynamic limit),
    producing correlations that span the entire system.
\end{enumerate}

\subsection{Mechanism of Self-Organization}

The interplay between slow driving (tree growth at rate $p$) and fast relaxation
(fire spread) creates a feedback loop:
\begin{itemize}
    \item If tree density is below the percolation threshold, fires remain small,
    and growth steadily increases connectivity.
    \item If density exceeds the threshold, a single lightning strike triggers
    a system-spanning fire, abruptly reducing density.
    \item The stationary state—the SOC state—is the percolation threshold itself.
\end{itemize}

This is analogous to the Bak-Tang-Wiesenfeld sandpile: grains are added slowly, and
avalanches of all sizes relax the system, keeping it perpetually at the angle of
repose \cite{bak1987self}.

\subsection{Finite-Size Effects}

With $L = 50$, the maximum avalanche size is $2500$ cells. Finite-size effects
truncate the power-law tail—the true power law would extend to infinity only in the
thermodynamic limit $L \to \infty$. A scaling analysis with varying $L$ would reveal
the finite-size scaling function and allow extraction of the exponent $\tau$.

% ==============================================================
\section{Conclusion}
\label{sec:conclusion}

This work has demonstrated that the forest fire model is a compelling example of
Self-Organized Criticality. The system exhibits all defining features: stochastic
noise (tree growth and lightning), cascading avalanches (fire spread), and a
connectivity-dependent critical state (percolation threshold). The Python simulation
on a $50 \times 50$ grid confirms a power-law-like avalanche size distribution,
consistent with scale-free dynamics.

Key findings are:
\begin{itemize}
    \item The avalanche size distribution follows $N(S) \sim S^{-\tau}$, indicative
    of criticality.
    \item Tree density (connectivity) determines the regime: sparse forests produce
    exponential distributions, while critical density yields power laws.
    \item The system self-organizes to the critical state without any external
    parameter tuning, through the feedback between slow growth and rapid fire
    relaxation.
\end{itemize}

Future work could include: (i) larger grids to reduce finite-size effects; (ii)
extraction of the scaling exponent $\tau$ via linear regression on the log-log plot;
(iii) spatial analysis of fire cluster shapes; and (iv) comparison with empirical
wildfire data to assess ecological validity.

% ==============================================================
\section*{Acknowledgements}

The author acknowledges the use of the following AI assistance tools in completing
this project:
\begin{itemize}
    \item \textbf{Claude (Anthropic)} — for generating the project report structure,
    LaTeX formatting, theoretical exposition, and analysis narrative based on the
    provided simulation code and assignment requirements.
    \item \textbf{ChatGPT (OpenAI)} — used for debugging and resolving minor issues in the simulation code.  
    \item \textbf{Gemini (Google)} — used for supplementary research and background understanding of concepts related to complexity science and self-organized criticality.  
 
\end{itemize}

All intellectual content has been reviewed and verified by the author. The simulation
code was written independently by the author.

% ==============================================================
\section*{List of Prompts Used}

The following prompts were submitted to Claude (Anthropic) to assist in writing this
report:

\begin{enumerate}
    \item \textit{``Explain self-organized criticality and forest fire models for conceptual clarity.'}
    \item \textit{``Suggest improvements for formatting a LaTeX report.'}
    \item \textit{``Help debug this Python code for a forest fire simulation.'}
\end{enumerate}

% ==============================================================
\bibliographystyle{unsrt}
\begin{thebibliography}{9}

\bibitem{bak1987self}
P. Bak, C. Tang, and K. Wiesenfeld,
\textit{Self-organized criticality: An explanation of the 1/f noise},
Physical Review Letters, \textbf{59}(4), 381--384, 1987.

\bibitem{drossel1992self}
B. Drossel and F. Schwabl,
\textit{Self-organized critical forest-fire model},
Physical Review Letters, \textbf{69}(11), 1629--1632, 1992.

\bibitem{bak1996nature}
P. Bak,
\textit{How Nature Works: The Science of Self-Organized Criticality},
Copernicus Press, New York, 1996.

\bibitem{jensen1998self}
H. J. Jensen,
\textit{Self-Organized Criticality: Emergent Complex Behavior in Physical and
Biological Systems},
Cambridge University Press, Cambridge, 1998.

\bibitem{christensen2005complexity}
K. Christensen and N. R. Moloney,
\textit{Complexity and Criticality},
Imperial College Press, London, 2005.

\bibitem{stauffer1992introduction}
D. Stauffer and A. Aharony,
\textit{Introduction to Percolation Theory}, 2nd ed.,
Taylor \& Francis, London, 1992.

\end{thebibliography}

\end{document}
